\documentclass{article}
\usepackage[utf8]{inputenc}
\usepackage[spanish]{babel}
\usepackage{geometry}
\usepackage{listings}
\usepackage{xcolor}
\geometry{a4paper, margin=1in}

\title{Propuesta de Estructura de Base de Datos: Meta-Pathfinder}
\author{Arquitectura de Datos}
\date{Mayo 2026}

\begin{document}

\maketitle

\section{Introducción}
Esta propuesta define la arquitectura relacional (o no relacional optimizada) para el sistema \textbf{Meta-Pathfinder}. La estructura está diseñada para soportar la persistencia de usuarios, el seguimiento de eventos cognitivos en tiempo real y la analítica metacognitiva necesaria para el motor de Inteligencia Artificial.

\section{Esquema de Colecciones/Tablas}

\subsection{Entidad: Users (Usuarios)}
Almacena la información de identidad y perfil de los estudiantes y administradores.
\begin{itemize}
    \item \texttt{id}: UUID / String (Primary Key)
    \item \texttt{name}: String (Nombre)
    \item \texttt{last\_name}: String (Apellido)
    \item \texttt{email}: String (Identificador único de acceso)
    \item \texttt{password\_hash}: String (Almacenamiento seguro)
    \item \texttt{role}: Enum [``student'', ``admin'']
    \item \texttt{avatar\_url}: String (Derivado de iniciales o personalizado)
    \item \texttt{created\_at}: Timestamp
\end{itemize}

\subsection{Entidad: Evaluations (Evaluaciones)}
Registra los resultados de las pruebas técnicas y metacognitivas.
\begin{itemize}
    \item \texttt{id}: UUID
    \item \texttt{user\_id}: UUID (Foreign Key $\to$ Users.id)
    \item \texttt{type}: Enum [``pretest'', ``challenge'', ``desfase'']
    \item \texttt{theme}: String (ej. ``Lógica'', ``Ciclos'', ``Funciones'')
    \item \texttt{score}: Float (0-100)
    \item \texttt{answers}: JSON (Detalle de las respuestas dadas)
    \item \texttt{created\_at}: Timestamp
\end{itemize}

\subsection{Entidad: CognitiveEvents (Eventos de Seguimiento)}
Captura el micro-comportamiento dinámico durante las evaluaciones.
\begin{itemize}
    \item \texttt{id}: UUID
    \item \texttt{user\_id}: UUID (Foreign Key)
    \item \texttt{type}: String (ej. ``PHASE\_START'', ``IMPULSE\_DETECTED'', ``SCROLL\_PATTERN'')
    \item \texttt{metadata}: JSON (ej. \{ ``latency'': 4.5, ``tab\_switches'': 2 \})
    \item \texttt{timestamp}: Timestamp
\end{itemize}

\subsection{Entidad: MetacognitiveProfile (Perfil de Estado)}
Mantiene el estado actual de las métricas del estudiante para personalización rápida.
\begin{itemize}
    \item \texttt{user\_id}: UUID (PK/FK)
    \item \texttt{calibration}: Float (Índice de precisión de autopercepción)
    \item \texttt{cognitive\_load\_avg}: Float (Histórico de carga cognitiva)
    \item \texttt{transfer\_readiness}: Float (Disponibilidad de transferencia)
    \item \texttt{latent\_strategies}: Array$<$String$>$ (Estrategias detectadas por la IA)
    \item \texttt{last\_updated}: Timestamp
\end{itemize}

\section{Consideraciones de Integración (Firebase/NoSQL)}
Si se utiliza una base de datos NoSQL como Firestore, se recomienda la siguiente jerarquía de documentos:
\begin{lstlisting}[basicstyle=\small\ttfamily, breaklines=true]
/users/{userId}
  /evaluations/{evalId}
  /cognitive_log/{eventId}
  /metacognition/data (documento unico por usuario)
\end{lstlisting}

\end{document}
